\documentclass{article}
\usepackage{graphicx}
\usepackage{amsmath}
\usepackage{amsfonts}
\usepackage{bookmark}
\usepackage{booktabs}
\usepackage{exercises}
\graphicspath{ {} }
\usepackage{indentfirst}
\usepackage{amsthm}
\usepackage{yhmath}
\usepackage{geometry}
\usepackage{placeins}
\usepackage{csquotes}
\usepackage{pgfplots}
\newcommand*{\oldneg}{\mathord{\sim}}
\newcommand\showdiv[1]{\overline{\smash{)}#1}}
\newcommand\Mydiv[1]{%
\smash{\raise.3ex\hbox{$\big)$}}\mkern-8mu
        \overline{\enspace\strut#1}}
\setlength{\tabcolsep}{5pt}
\theoremstyle{definition}
\newtheorem{definition}{Definition}

\setlength{\parindent}{.5in}
\setlength{\parskip}{12pt}
\renewcommand{\arraystretch}{2}
\geometry{margin=1in}
\newcommand{\?}{\stackrel{?}{=}}
\pgfplotsset{width=7cm,compat=1.18}\usepgfplotslibrary{statistics}

\title{9.2 Answers}
\begin{document}
\maketitle

\section*{1}
An angle is in standard position if its vertex is at the origin and its initial side lies on the positive $x$-axis.

\section*{2}
A negative angle is clockwise, while a positive angle is counter-clockwise.

\section*{3}
B

\section*{4}
C

\section*{5}
A

\section*{14}
C

\section*{17}
All of the coterminal angles of $-125^\circ$ are $-125+360n$ where $n$ is an integer. One negative example is $-125-360=-485^\circ$, and a positive example is
$-125+360=235^\circ$.

\section*{18}
All of the coterminal angles of $820^\circ$ are $820+360n$ where $n$ is an integer. One negative example is $820-360(3)=-260^\circ$ and one positive example is
$820-360=460^\circ$.

\section*{21}
All of the coterminal angles of $\frac{28\pi}{9}$ are $\frac{28\pi}{9}+2\pi n$ where $n$ is an integer. One negative example is
\[ \frac{28\pi}{9}-4\pi=\frac{28\pi-36\pi}{9}=-\frac{8\pi}{9} \]
and one positive example is
\[ \frac{28\pi}{9}-2\pi=\frac{28\pi-18\pi}{9}=\frac{10\pi}{9} \]

\section*{22}
All of the coterminal angles of $\frac{20\pi}{3}$ are $\frac{20\pi}{3}+2\pi n$ where $n$ is an integer. One negative example is
\[ \frac{20\pi}{3}-8\pi=\frac{20\pi-24\pi}{3}=-\frac{4\pi}{3} \]
and one positive example is
\[ \frac{20\pi}{3}+2\pi=\frac{20\pi+6\pi}{3}=\frac{26\pi}{3} \]

\section*{24}
\[ 315^\circ=315\times\frac{\pi}{180}=\frac{7\pi}{4} \]

\section*{25}
\[ -260^\circ=-260\times\frac{\pi}{180}=-\frac{13\pi}{9} \]

\section*{26}
\[ 500^\circ=500\times\frac{\pi}{180}=\frac{25\pi}{9} \]

\section*{28}
\[ -\frac{\pi}{4}=-\frac{\pi}{4}\times\frac{180}{\pi}=-45^\circ \]

\section*{29}
\[ 5\pi=5\pi\times\frac{180}{\pi}=900^\circ \]

\section*{30}
\[ \frac{14\pi}{15}=\frac{14\pi}{15}\times\frac{180}{\pi}=168^\circ \]

\section*{31}
\[ \frac{13\pi}{6}=\frac{13\pi}{6}\times\frac{180}{\pi}=390^\circ \] Thus the answer is B.

\section*{32}
We know that the arc length is $\frac{\pi}{6}\times 4=\frac{2\pi}{3}$ inches and the area of the sector is \[ \frac{1}{2}\times 4^2\times\frac{\pi}{6}=\frac{4\pi}{3} \] square inches.

\section*{33}
We know that the arc length is $\frac{5\pi}{12}\times 3=\frac{5\pi}{4}$ meters and the area of the sector is
\[ \frac{1}{2}\times 3^2\times\frac{5\pi}{12}=\frac{45\pi}{24}=\frac{15\pi}{8} \]

\section*{35}
We know that the central angle measured in radians is $150\times\frac{\pi}{180}=\frac{5\pi}{6}$, meaning the arc length is 
$\frac{5\pi}{6}\times 12=10\pi$ ft. and the area of the sector is \[ \frac{1}{2}\times 12^2\times\frac{5\pi}{6}=60\pi \] square feet.

\section*{36}
We know that the central angle measured in radians is $25\times\frac{\pi}{180}=\frac{5\pi}{36}$, meaning the arc length is
$\frac{5\pi}{36}\times 18=\frac{5\pi}{2}$ meters and the area of the sector is
\[ \frac{1}{2}\times 18^2\times\frac{5\pi}{36}=\frac{45\pi}{2} \]

\section*{38}
The error here is that the formula for finding the area of a sector does not work with an angle measured in degrees, but only with an angle measured in radians. 
In the problem, we know that the central angle measured in radians is $40\times\frac{\pi}{180}=\frac{2\pi}{9}$, meaning the area of the sector is
\[ \frac{1}{2}\times 6^2\times\frac{2\pi}{9}=4\pi \] square centimeters.

\section*{39}
We know that $\frac{\pi}{3}=\frac{\pi}{3}\times\frac{180}{\pi}=60^\circ$. In a 30-60 right triangle, we know that
$\cos 60=\frac{1}{2}$. Thus $\cos\frac{\pi}{3}=\frac{1}{2}$.

\section*{40}
We know that $\frac{\pi}{4}=\frac{\pi}{4}\times\frac{180}{\pi}=45^\circ$. In a 45-45 right triangle, we know that
$\sin 45=\frac{1}{\sqrt{2}}=\frac{\sqrt{2}}{2}$. Thus $\sin\frac{\pi}{4}=\frac{\sqrt{2}}{2}$.

\section*{41}
We know that $\frac{\pi}{6}=\frac{\pi}{6}\times\frac{180}{\pi}=30^\circ$. In a 30-60 right triangle, we know that
$\tan 30=\frac{1}{\sqrt{3}}=\frac{\sqrt{3}}{3}$. Thus $\tan\frac{\pi}{6}=\frac{\sqrt{3}}{3}$.

\section*{44}
We know that $\frac{\pi}{6}=30^\circ$, and because 
$\cos 30=\frac{\sqrt{3}}{2}$, we know that $\cos\frac{\pi}{6}=\frac{\sqrt{3}}{2}$.

\section*{47}
I was not understanding the example of the notation given in the problem.

\section*{51}
(a): We are given that they spun 15 revolutions per minute, meaning the spun $\frac{15}{60}=\frac{1}{4}$ revolutions per second. Because there are $2\pi$ radians in a circle, we know that 
they spun $\frac{2\pi}{4}=\frac{\pi}{2}$ radians each second.


(b): We know that the sector has a central angle of $\frac{\pi}{2}$ and a radius of 29 feet (because the diameter is 58 feet), meaning the length of the arc is $29\times\frac{\pi}{2}=\frac{29\pi}{2}$ feet.

\section*{52}
(a): We know that the radius is 21.16 meters and the central angle is $34.92^\circ$. Thus the central angle, in radians, is
$34.92\times\frac{\pi}{180}=0.194\pi$. Thus the length of the arc is $21.16\times 0.194\pi=4.10504\pi$ meters.


(b): The area of the arc is \[ \frac{1}{2}\times 21.16^2\times 0.194\pi\approx 43.43\pi \] square meters.

\section*{53}
(a): The length of the arc is $42\times\frac{\pi}{8}=\frac{21\pi}{4}$ inches.


(b): We know that, for one stair, the angle turned is $\frac{\pi}{8}$. Thus for 16 turns, the angle is $16\times\frac{\pi}{8}=2\pi$.


(c): The area of the sector for just one step is $\frac{1}{2}\times 42^2\times\frac{\pi}{8}=\frac{441\pi}{4}$ square inches. Thus, the total area for the 15 steps is
\[ 15\times\frac{441\pi}{4}=\frac{6615\pi}{4} \] square inches.

\section*{54}
(a): We know that the radius of the entire sector is $\frac{53}{8}$ inches. Because there are twenty such sectors in the circle, we know that each has a central angle of 
$\frac{2\pi}{20}=\frac{\pi}{10}$ radians. Thus the area of the entire sector is $\frac{1}{2}\times\frac{53}{8}^2\times\frac{\pi}{10}=\frac{2809\pi}{1280}$. The area of the double area is 
equal to the area of the entire sector minus the area of the sector with the same central angle and a radius of $\frac{50}{8}$ inches. The area of this sector is
\[ \frac{1}{2}\times\frac{50}{8}^2\times\frac{\pi}{10}=\frac{625\pi}{320} \] meaning the area of the double sector is
\[ \frac{2809\pi}{1280}-\frac{625\pi}{320}=\frac{2809\pi-2500\pi}{1280}=\frac{309\pi}{1280} \] square inches. The area of triple region is equal to the area of the sector with a radius of 
$\frac{15}{4}+\frac{3}{8}=\frac{33}{8}$ inches minus the area of the sector with a radius of $\frac{15}{4}$ inches. The area of the first sector is
\[ \frac{1}{2}\times\frac{33}{8}^2\times\frac{\pi}{10}=\frac{1089\pi}{1280} \] and the area of the second sector is
\[ \frac{1}{2}\times\frac{15}{4}^2\times\frac{\pi}{10}=\frac{45\pi}{64} \] meaning the area of the triple area is
\[ \frac{1089\pi}{1280}-\frac{45\pi}{64}=\frac{1089\pi-900\pi}{1280}=\frac{189\pi}{1280} \] 


(b): The probability that the dart lands in the triple region is equal to the ratio of the area of the triple sector to the entire sector, which is
\[ \dfrac{\frac{189\pi}{1280}}{\frac{2809\pi}{1280}}=\frac{189}{2809}\approx 0.067 \]
The probability that the dart lands in the double region is
\[ \dfrac{\frac{309\pi}{1280}}{\frac{2809\pi}{1280}}=\frac{309}{2809}\approx 0.110 \]

\end{document}
